% Options for packages loaded elsewhere
\PassOptionsToPackage{unicode}{hyperref}
\PassOptionsToPackage{hyphens}{url}
\documentclass[
  ignorenonframetext,
]{beamer}
\newif\ifbibliography
\usepackage{pgfpages}
\setbeamertemplate{caption}[numbered]
\setbeamertemplate{caption label separator}{: }
\setbeamercolor{caption name}{fg=normal text.fg}
\beamertemplatenavigationsymbolsempty
% remove section numbering
\setbeamertemplate{part page}{
  \centering
  \begin{beamercolorbox}[sep=16pt,center]{part title}
    \usebeamerfont{part title}\insertpart\par
  \end{beamercolorbox}
}
\setbeamertemplate{section page}{
  \centering
  \begin{beamercolorbox}[sep=12pt,center]{section title}
    \usebeamerfont{section title}\insertsection\par
  \end{beamercolorbox}
}
\setbeamertemplate{subsection page}{
  \centering
  \begin{beamercolorbox}[sep=8pt,center]{subsection title}
    \usebeamerfont{subsection title}\insertsubsection\par
  \end{beamercolorbox}
}
% Prevent slide breaks in the middle of a paragraph
\widowpenalties 1 10000
\raggedbottom
\AtBeginPart{
  \frame{\partpage}
}
\AtBeginSection{
  \ifbibliography
  \else
    \frame{\sectionpage}
  \fi
}
\AtBeginSubsection{
  \frame{\subsectionpage}
}
\usepackage{iftex}
\ifPDFTeX
  \usepackage[T1]{fontenc}
  \usepackage[utf8]{inputenc}
  \usepackage{textcomp} % provide euro and other symbols
\else % if luatex or xetex
  \usepackage{unicode-math} % this also loads fontspec
  \defaultfontfeatures{Scale=MatchLowercase}
  \defaultfontfeatures[\rmfamily]{Ligatures=TeX,Scale=1}
\fi
\usepackage{lmodern}
\ifPDFTeX\else
  % xetex/luatex font selection
\fi
% Use upquote if available, for straight quotes in verbatim environments
\IfFileExists{upquote.sty}{\usepackage{upquote}}{}
\IfFileExists{microtype.sty}{% use microtype if available
  \usepackage[]{microtype}
  \UseMicrotypeSet[protrusion]{basicmath} % disable protrusion for tt fonts
}{}
\makeatletter
\@ifundefined{KOMAClassName}{% if non-KOMA class
  \IfFileExists{parskip.sty}{%
    \usepackage{parskip}
  }{% else
    \setlength{\parindent}{0pt}
    \setlength{\parskip}{6pt plus 2pt minus 1pt}}
}{% if KOMA class
  \KOMAoptions{parskip=half}}
\makeatother
\setlength{\emergencystretch}{3em} % prevent overfull lines
\providecommand{\tightlist}{%
  \setlength{\itemsep}{0pt}\setlength{\parskip}{0pt}}
\usepackage{bookmark}
\IfFileExists{xurl.sty}{\usepackage{xurl}}{} % add URL line breaks if available
\urlstyle{same}
\hypersetup{
  hidelinks,
  pdfcreator={LaTeX via pandoc}}

\author{\texorpdfstring{}{}}
\date{}

\begin{document}

\section{Introduction}\label{introduction}

\begin{frame}{The Reproducibility Crisis and Its Structural Roots}
\protect\phantomsection\label{the-reproducibility-crisis-and-its-structural-roots}
The ``reproducibility crisis'' in computational and biological sciences
is not merely a cultural failure but a structural one. When research
papers are divorced from the code that generated their figures, when
datasets are shared as unlabeled archives, and when build processes
exist only as tribal knowledge in a graduate student's shell history,
the conditions for irreproducibility are baked into the workflow itself.
A 2016 \emph{Nature} survey of 1,576 researchers found that 70\% had
tried and failed to reproduce another scientist's experiments, and more
than half had failed to reproduce their own
{[}@baker2016reproducibility{]}. The economic cost is staggering:
Freedman et al.~estimate that the biomedical industry alone loses \$28
billion annually to irreproducible preclinical research
{[}@freedman2015economics{]}.

The root cause is fragmentation. A typical research project scatters its
artifacts across disconnected tools: LaTeX editors for prose, Jupyter
notebooks for analysis, ad-hoc shell scripts for figure generation, and
manual copy-paste for integrating results into manuscripts. Each
boundary between tools is a potential locus of desynchronization. The
version of the figure in the PDF may not match the version of the code
that ostensibly generated it. The test suite, if it exists at all,
likely tests the code in isolation from the rendering pipeline. Peng
{[}@peng2011reproducible{]} argues that reproducibility in computational
science requires, at minimum, that the data and code underlying a
published result be available for independent verification---yet the
tools for enforcing this standard remain ad hoc. Indeed, even the
terminology is fractured: Barba {[}@barba2018terminologies{]} documents
how ``reproducibility,'' ``replicability,'' and ``repeatability'' carry
conflicting definitions across disciplines, undermining cross-field
standards.
\end{frame}

\begin{frame}{Related Work}
\protect\phantomsection\label{related-work}
Gentleman and Temple Lang {[}@gentleman2007research{]} introduced the
concept of a \emph{research compendium}---a single unit of scholarly
communication bundling code, data, and narrative. This vision has driven
two decades of tooling, which can be grouped into four categories:
workflow managers, literate programming systems, containerization
approaches, and best-practice frameworks.

\begin{block}{Workflow Managers}
\protect\phantomsection\label{workflow-managers}
\textbf{Snakemake} {[}@koster2012snakemake{]} uses a rule-based,
Python-derived DSL to specify computational workflows as directed
acyclic graphs of file-producing steps. It supports containerized
execution via Conda and Singularity environments. However, Snakemake
focuses exclusively on computational pipeline orchestration and does not
integrate manuscript rendering, testing enforcement, or provenance
watermarking.

\textbf{Nextflow} {[}@ditommaso2017nextflow{]} employs a dataflow
programming paradigm with native support for container-based execution
across heterogeneous computing environments (local, SLURM, AWS). Like
Snakemake, Nextflow excels at bioinformatics pipeline parallelism but
does not address manuscript production, document integrity, or the
testing--publication coupling that characterizes research
reproducibility.

\textbf{CWL} (Common Workflow Language) {[}@amstutz2016cwl{]} provides a
portable, YAML-based standard for describing computational workflows and
their dependencies. Its strength lies in interoperability across
execution engines (cwltool, Toil, Arvados), but it requires external
tooling for manuscript generation and offers no built-in testing or
provenance framework.
\end{block}

\begin{block}{Literate Programming and Publication Tools}
\protect\phantomsection\label{literate-programming-and-publication-tools}
Knuth's literate programming {[}@knuth1984literate{]} established the
principle that programs should be authored as documents intended for
human comprehension. Schulte et al. {[}@schulte2012multilanguage{]}
extended this to multi-language computing environments (Org-mode),
demonstrating that literate programming could span languages and output
formats.

\textbf{Quarto} {[}@allaire2024quarto{]} extends the R Markdown
tradition to support Python, Julia, and Observable, rendering to PDF,
HTML, and Word. Quarto integrates code execution with document
rendering, achieving a modern form of literate programming, but it does
not enforce testing thresholds, manage multi-project repositories, or
provide cryptographic provenance.

\textbf{Jupyter Book} {[}@kluyver2016jupyter{]} builds on Jupyter
notebooks to produce publication-quality documents via Sphinx. While
powerful for interactive exploration, Jupyter's notebook format
introduces execution-order fragility and does not naturally support the
separation of logic from orchestration that characterizes maintainable
research software.

\textbf{R Markdown} {[}@xie2015dynamic{]} pioneered knitr-based dynamic
documents that weave code and prose. Its ecosystem is rich but
R-centric, and it lacks the multi-project management, infrastructure
testing, and provenance embedding that characterize the Docxology
Template.
\end{block}

\begin{block}{Containerization and Environment Capture}
\protect\phantomsection\label{containerization-and-environment-capture}
\textbf{Docker} {[}@boettiger2015docker{]} addresses reproducibility at
the environment level---packaging operating system, libraries, and code
into portable containers. While Docker solves the ``works on my
machine'' problem, containerization is complementary to, not a
replacement for, the architectural concerns addressed here: Docker does
not enforce testing, embed provenance, or manage multi-project
manuscript workflows.
\end{block}

\begin{block}{Best-Practice Frameworks and Data Standards}
\protect\phantomsection\label{best-practice-frameworks-and-data-standards}
Wilson et al. {[}@wilson2017good{]} define ``good enough'' practices for
scientific computing, emphasizing version control, testing, and
documentation. Sandve et al. {[}@sandve2013ten{]} propose ten rules for
reproducible computational research. Stodden et al.
{[}@stodden2016enhancing{]} advocate for enhanced computational method
transparency. The FAIR principles {[}@wilkinson2016fair{]}---Findable,
Accessible, Interoperable, Reusable---establish a standard for data
stewardship that has been widely adopted by funding agencies and
journals. Barker et al. {[}@barker2022fair4rs{]} extend these principles
specifically to research software via the FAIR4RS initiative,
recognizing that software has distinct requirements---executability,
composability, and dependency management---that data-centric FAIR does
not address. Garijo et al. {[}@garijo2024fairsoft{]} operationalize
FAIR4RS through the FAIRsoft evaluator, an automated assessment
framework that scores research software against 17+ quality indicators
including executability, metadata richness, and documentation
completeness. Nüst et al. {[}@nust2017containerization{]} introduce the
\emph{executable research compendium} (ERC), extending Gentleman and
Temple Lang's compendium concept with containerized, interactive
reproduction environments. The W3C PROV data model
{[}@moreau2013provdm{]} provides a formal vocabulary for expressing
provenance records, while in-toto {[}@torresarias2019intoto{]} provides
a framework for end-to-end software supply chain integrity verification.
These frameworks articulate \emph{what} reproducible research requires
but do not provide integrated \emph{how}---they lack the tooling,
enforcement mechanisms, and architectural patterns that translate
standards into practice.
\end{block}

\begin{block}{The Gap}
\protect\phantomsection\label{the-gap}
Despite advances in FAIR4RS principles {[}@barker2022fair4rs{]} and
automated FAIR software assessment {[}@garijo2024fairsoft{]}, no
existing system integrates all six concerns into a single enforced
pipeline: (1) end-to-end pipeline orchestration with testing
enforcement, (2) multi-format manuscript rendering, (3) cryptographic
provenance embedding, (4) multi-project repository management, (5)
FAIR-aligned software stewardship, and (6) AI-agent collaboration via
structured documentation. FAIR4RS provides the vocabulary; the Docxology
Template provides the enforcement mechanism.
\end{block}
\end{frame}

\begin{frame}[fragile]{The Docxology Template: An Integrated Solution}
\protect\phantomsection\label{the-docxology-template-an-integrated-solution}
The \emph{Docxology Template} was conceived as a structural antidote to
this fragmentation. Rather than adding reproducibility as an
afterthought---a Docker container wrapping an already-disjointed
workflow {[}@boettiger2015docker{]}---the template enforces integrity at
the architectural level. It realizes Gentleman and Temple Lang's
research compendium vision {[}@gentleman2007research{]} at repository
scale, bundling code, data, tests, manuscripts, and provenance into a
single, pipeline-enforced system. It stands on four primary pillars:

\begin{enumerate}
\item
  \textbf{Ergonomic Modularity}: A Two-Layer Architecture cleanly
  separates globally shared infrastructure (logging, rendering,
  validation, steganography) from project-specific logic (manuscripts,
  scripts, data). Ten infrastructure subpackages comprising 137 Python
  modules provide reusable services; projects consume them without
  modification.
\item
  \textbf{Execution Integrity}: A Zero-Mock testing policy where
  pipeline advancement is contingent on test passage. Infrastructure
  tests must achieve 60\% coverage; project tests must achieve 90\%. All
  tests use real filesystem operations, real subprocess calls, and real
  network connections---no mock objects, no fake services, no synthetic
  test doubles. Over 3,000 infrastructure tests and 350+ project tests
  enforce this standard.
\item
  \textbf{Automated Provenance}: Steganographic watermarking and
  cryptographic hashing are integrated directly into the rendering
  pipeline. Every generated PDF carries a SHA-256 fingerprint, an
  alpha-channel text overlay encoding the build timestamp and commit
  hash, and optionally a QR code linking to the repository. Provenance
  is not asserted by policy; it is enforced by architecture.
\item
  \textbf{AI-Agent Collaboration}: A Documentation Duality standard
  equips every directory with both \texttt{README.md} (for human
  researchers) and \texttt{AGENTS.md} (for AI coding assistants).
  Supplementary \texttt{CLAUDE.md} files provide system-level
  instructions. This enables a novel human--AI pair-programming workflow
  where AI agents can navigate, modify, and extend the codebase with
  full architectural awareness.
\end{enumerate}
\end{frame}

\begin{frame}[fragile]{Scope and Contributions}
\protect\phantomsection\label{scope-and-contributions}
This paper is itself a product of the template it describes---a
self-referential demonstration of the system's capabilities. Our
contributions are:

\begin{itemize}
\tightlist
\item
  A formal description of the Two-Layer Architecture and Standalone
  Project Paradigm that enables N independent research projects to share
  infrastructure without coupling.
\item
  A detailed specification of the eight-stage build pipeline, from
  environment sanitization through LLM-assisted review.
\item
  A comparative analysis positioning the template against Snakemake,
  Nextflow, CWL, Quarto, and Jupyter Book.
\item
  An empirical evaluation of the system across three heterogeneous
  exemplar projects, demonstrating scalability, coverage metrics, and
  pipeline timing.
\item
  A security analysis of the steganographic provenance layer, including
  threat models and tamper-detection capabilities.
\item
  An open-source reference implementation available at
  \texttt{github.com/docxology/template}.
\end{itemize}
\end{frame}

\begin{frame}{Paper Organization}
\protect\phantomsection\label{paper-organization}
The Methods (§2) describe the Two-Layer Architecture, Thin Orchestrator
pattern, pipeline stages, and AI collaboration model. Results (§3)
present quantitative metrics from multi-project execution, coverage
analysis, and steganographic benchmarks. The Discussion (§4) addresses
implications, limitations, and future directions. The Infrastructure
Module Reference (§5) provides detailed documentation for all ten
subpackages. Security and Provenance (§6) describes the steganographic
and cryptographic integrity layer. Appendices (§7) provide pipeline,
configuration, and comparative references.
\end{frame}

\end{document}
